\documentclass[../../dis.tex]{subfiles}
\begin{document}
\chapter{L2 - Relational Model \& SQL Basics}\label{chap:l2}
\vfill
\section{Relational model}
The relational model is the foundation of modern database systems. It represents all data as tables and provides a clean, mathematically grounded framework for querying and constraining data.

\subsection{Core concepts}
A \textbf{database} is a set of named \recall{def:relation}{relations}{named tables with rows and columns}. Each attribute carries a \textbf{type} (domain): integer, real, string, file formats, enumerated, etc.

\begin{concept}{Tuple}
A single row in a relation, containing one value for each attribute. All tuples in a relation are distinct and have no inherent order.
\end{concept}

Running example:

\begin{minipage}[t]{0.48\textwidth}
\begin{dbtable}{Students}{lllcc}
\textbf{sid} & \textbf{name} & \textbf{login} & \textbf{age} & \textbf{gpa} \\ \hline
50000 & Dave & dave@cs & 19 & 3.3 \\ \hline
53666 & Jones & jones@cs & 18 & 3.4 \\ \hline
53688 & Smith & smit@ee & 18 & 3.2 \\ \hline
\ldots & \ldots & \ldots & \ldots & \ldots \\
\end{dbtable}
\end{minipage}
\hfill
\begin{minipage}[t]{0.48\textwidth}
\begin{dbtable}{Colleges}{llc}
\textbf{name} & \textbf{location} & \textbf{strength} \\ \hline
MIT & USA & 10000 \\ \hline
Oxford & UK & 22000 \\ \hline
EPFL & CH & 9000 \\ \hline
\ldots & \ldots & \ldots \\
\end{dbtable}
\end{minipage}

\subsection{Schema vs.\ instance}
A relation has two complementary aspects: its \recall{def:schema}{schema}{structural definition with attribute names and types} and its content at a given moment.

Example schema: \texttt{Students(\underline{sid}: string, name: string, login: string, age: integer, gpa: real)}.

\begin{concept}{Instance}
The actual set of tuples stored in a relation at a given point in time.
\end{concept}

Two measures describe an instance:
\begin{dashlist}
    \item \textbf{Cardinality}: the number of tuples (rows).
    \item \textbf{Arity} (or degree): the number of attributes (columns).
\end{dashlist}

\subsection{Null values}
Not every attribute value is always known. A student who has not received any grades yet has no meaningful GPA.

\begin{concept}{Null value}
A special marker representing an ``unknown'' or ``undefined'' attribute value. Nulls are not the same as zero or an empty string.
\end{concept}

\section{Keys and integrity constraints}
\recall{def:key}{Keys}{attributes that uniquely identify tuples} enforce uniqueness and link relations together. Here we refine the hierarchy of key types and formalize integrity constraints.

\subsection{Key hierarchy}

\begin{concept}{Superkey}
A set of attributes such that no two distinct tuples can share the same values in all those fields. Any superset of a key is also a superkey.
\end{concept}

\begin{concept}{Candidate key}
A \emph{minimal} superkey: it is a superkey, and no proper subset of its fields is also a superkey.
\end{concept}

\begin{concept}{Primary key}
The candidate key chosen by the database administrator (DBA) as \emph{the} key of the relation.
\end{concept}

Example with \textbf{Students}(\underline{sid}, name, login, age, gpa): both \texttt{sid} and \texttt{login} are candidate keys. The DBA picks \texttt{sid} as the primary key.

\begin{dashlist}
    \item Every superset of \texttt{sid} (e.g., \texttt{\{sid, name\}}, \texttt{\{sid, login, gpa\}}, \ldots) is a superkey but not a candidate key.
    \item \texttt{sid} alone and \texttt{login} alone are the only candidate keys.
\end{dashlist}

\subsection{Integrity constraints}
A key is an \textbf{integrity constraint} (IC): a condition that must hold on every valid database instance.

\recall{def:foreignkey}{Foreign keys}{attributes referencing the primary key of another relation} establish \textbf{referential integrity}: every foreign-key value must match an existing key in the referenced relation.

\newpage
\section{SQL: Structured Query Language}
SQL is the standard language for interacting with relational databases. It splits into two sub-languages.

\subsection{DDL and DML}
\begin{concept}{Data Definition Language (DDL)}
The subset of SQL used to create, modify, and delete relations, specify constraints, and administer users and security.
\end{concept}

\begin{concept}{Data Manipulation Language (DML)}
The subset of SQL used to query, insert, update, and delete tuples.
\end{concept}

\subsection{Most common SQL commands}
\begin{sql}
    CREATE TABLE <name> (<field> <domain>, ...)
    INSERT INTO <name> (<field names>) VALUES (<field values>)
    DELETE FROM <name> WHERE <condition>
    UPDATE <name> SET <field name> = <value> WHERE <condition>
    SELECT <fields> FROM <name> WHERE <condition>
\end{sql}

\subsection{Creating relations}
A \lstinline[style=SQLStyle]{CREATE TABLE} statement defines a new relation by listing each attribute and its domain. The type of each field is enforced by the DBMS whenever tuples are added or modified.

\begin{minipage}[t]{0.48\textwidth}
\textbf{Students} -- holds data about students:
\begin{sql}
    CREATE TABLE Students
      (sid   CHAR(20),
       name  CHAR(20),
       login CHAR(10),
       age   INTEGER,
       gpa   FLOAT)
\end{sql}
\end{minipage}
\hfill
\begin{minipage}[t]{0.48\textwidth}
\textbf{Enrolled} -- holds course enrollments:
\begin{sql}
    CREATE TABLE Enrolled
      (sid   CHAR(20),
       cid   CHAR(20),
       grade CHAR(2))
\end{sql}
\end{minipage}

\subsection{Adding and deleting tuples}
Insert a single tuple:
\begin{sql}
    INSERT INTO Students (sid, name, login, age, gpa)
    VALUES ('53688', 'Smith', 'smith@cs', 18, 3.2)
\end{sql}

Delete all tuples satisfying a condition:
\begin{sql}
    DELETE FROM Students S WHERE S.name = 'Smith'
\end{sql}

\subsection{Keys in SQL}
The \lstinline[style=SQLStyle]{PRIMARY KEY} clause declares the primary key; \lstinline[style=SQLStyle]{UNIQUE} marks additional candidate keys.

\begin{minipage}[t]{0.48\textwidth}
\begin{sql}
    CREATE TABLE Students
      (sid   CHAR(20),
       name  CHAR(20),
       login CHAR(10),
       age   INTEGER,
       gpa   FLOAT,
       PRIMARY KEY(sid))
\end{sql}
\end{minipage}
\hfill
\begin{minipage}[t]{0.48\textwidth}
\begin{sql}
    CREATE TABLE Person
      (ssn      CHAR(9),
       name     CHAR(20),
       licence# CHAR(10),
       PRIMARY KEY(ssn),
       UNIQUE(licence#))
\end{sql}
\end{minipage}

\subsubsection*{Choosing keys carefully}
Keys must be used carefully: they encode real-world constraints. Consider \texttt{Enrolled}: ``for a given student and course, there is a single grade.''

\begin{minipage}[t]{0.48\textwidth}
\textbf{Correct} -- composite PK on \texttt{(sid, cid)}:
\begin{sql}
    CREATE TABLE Enrolled
      (sid   CHAR(20),
       cid   CHAR(20),
       grade CHAR(2),
       PRIMARY KEY(sid, cid))
\end{sql}
\end{minipage}
\hfill
\begin{minipage}[t]{0.48\textwidth}
\textbf{Wrong} -- PK on \texttt{sid} alone with \texttt{UNIQUE(cid, grade)} encodes a different constraint: ``students can take only one course, and no two students in a course receive the same grade.''
\end{minipage}

\subsection{Foreign keys and referential integrity}
\recall{def:foreignkey}{Foreign keys}{fields referencing the primary key of another relation} act as logical pointers between relations. If all foreign-key constraints are enforced, the database achieves \textbf{referential integrity} (no dangling references).

\begin{minipage}[t]{0.48\textwidth}
\begin{dbtable}{Students}{lllcc}
\textbf{sid} & \textbf{name} & \textbf{login} & \textbf{age} & \textbf{gpa} \\ \hline
50000 & Dave & dave@cs & 19 & 3.3 \\ \hline
53666 & Jones & jones@cs & 18 & 3.4 \\ \hline
53688 & Smith & smit@ee & 18 & 3.2 \\ \hline
\ldots & \ldots & \ldots & \ldots & \ldots \\
\end{dbtable}
\end{minipage}
\hfill
\begin{minipage}[t]{0.48\textwidth}
\begin{dbtable}{Enrolled}{llc}
\textbf{cid} & \textbf{sid} & \textbf{grade} \\ \hline
Carnatic101 & 53666 & C \\ \hline
Reggae203 & 50000 & B \\ \hline
Topology112 & 53666 & A \\ \hline
\ldots & \ldots & \ldots \\
\end{dbtable}
\end{minipage}

\smallskip
Here \texttt{Enrolled.sid} is a foreign key referencing \texttt{Students.sid}: every enrolled student must actually exist in \texttt{Students}.

\subsubsection*{Enforcing referential integrity}
The DBMS must decide what to do when a foreign-key constraint is threatened:
\begin{dashlist}
    \item \textbf{Insert} an Enrolled tuple with a non-existent \texttt{sid} $\rightarrow$ reject the insert.
    \item \textbf{Delete} a Students tuple that is referenced by Enrolled $\rightarrow$ several policies:
    \begin{dashlist}
        \item cascade-delete all Enrolled tuples that refer to it,
        \item disallow the deletion entirely,
        \item set the referencing \texttt{sid} to a default value,
        \item set it to \texttt{NULL} (denoting ``unknown'' or ``inapplicable'').
    \end{dashlist}
    \item \textbf{Update} the primary key of a Students tuple $\rightarrow$ same policies apply.
\end{dashlist}

\subsection{Integrity constraints}

\begin{concept}{Integrity constraint (IC)}
A condition that must be true for \emph{any} instance of the database (e.g., domain constraints, key constraints, foreign-key constraints). ICs are specified when the schema is defined and checked whenever relations are modified.
\end{concept}

\begin{concept}{Legal instance}
An instance of a relation that satisfies all specified ICs. The DBMS should not allow illegal instances.
\end{concept}

When the DBMS enforces ICs, stored data is more faithful to real-world meaning and data-entry errors are caught early.

\newpage
\section{Relational query languages}
\begin{concept}{Query language (QL)}
A language that allows manipulation and retrieval of data from a database. Unlike programming languages, QLs are not Turing-complete: they are designed for easy, efficient access to large data sets, not general computation.
\end{concept}

The relational model supports simple yet powerful QLs with a strong formal foundation based on logic, enabling significant optimization.

\subsection{Formal relational query languages}
Two mathematical query languages form the basis for real languages like SQL:

\begin{concept}{Relational Algebra}
An \emph{operational} language: it specifies \emph{how} to compute a result step by step. Useful for representing execution plans inside database engines.
\end{concept}

\begin{concept}{Relational Calculus}
A \emph{declarative} language: it specifies \emph{what} the result should look like, not how to compute it. Non-procedural.
\end{concept}

\begin{concept}{Codd's Theorem}
Relational algebra and relational calculus are equivalent in expressive power: for every query in one, there is an equivalent query in the other.
\end{concept}

Understanding both is key to understanding SQL and query processing.

\subsection{Preliminaries}
A query is applied to \emph{relation instances}, and the result of a query is also a relation instance.
\begin{dashlist}
    \item The \emph{schemas} of input relations are fixed (but the query runs over any legal instance).
    \item The schema of the \emph{result} is also fixed, determined by the query language constructs.
\end{dashlist}

Two notations are used:
\begin{dashlist}
    \item \textbf{Positional notation}: easier for formal definitions.
    \item \textbf{Named-field notation}: more readable; both are used in SQL.
\end{dashlist}

\subsection{Example schema and instances}
We use a university database throughout the relational algebra examples:

{\small
\begin{tabularx}{\textwidth}{@{}lX@{}}
\texttt{instructor} & \texttt{\underline{ID}, name, dept\_name, salary} \\
\texttt{course} & \texttt{\underline{course\_id}, title, dept\_name, credits} \\
\texttt{teaches} & \texttt{\underline{ID}, \underline{course\_id}, \underline{sec\_id}, \underline{semester}, \underline{year}} \\
\texttt{section} & \texttt{\underline{course\_id}, \underline{sec\_id}, \underline{semester}, \underline{year}, building, room\_number, time\_slot\_id} \\
\texttt{student} & \texttt{\underline{ID}, name, dept\_name, tot\_cred} \\
\texttt{takes} & \texttt{\underline{ID}, \underline{course\_id}, \underline{sec\_id}, semester, year, grade} \\
\end{tabularx}
}

Running example instance of \texttt{instructor}:

{\small
\begin{dbtable}{instructor}{clll}
\textbf{ID} & \textbf{name} & \textbf{dept\_name} & \textbf{salary} \\ \hline
22222 & Einstein & Physics & 95000 \\ \hline
12121 & Wu & Finance & 90000 \\ \hline
32343 & El Said & History & 60000 \\ \hline
45565 & Katz & Comp.\ Sci. & 75000 \\ \hline
98345 & Kim & Elec.\ Eng. & 80000 \\ \hline
76766 & Crick & Biology & 72000 \\ \hline
10101 & Srinivasan & Comp.\ Sci. & 65000 \\ \hline
58583 & Califieri & History & 62000 \\ \hline
83821 & Brandt & Comp.\ Sci. & 92000 \\ \hline
15151 & Mozart & Music & 40000 \\ \hline
33456 & Gold & Physics & 87000 \\ \hline
76543 & Singh & Finance & 80000 \\
\end{dbtable}
}
\subsection{Five basic operations}
Relational algebra has five primitive operators, split into two groups:

\textbf{Unary operators} (one input relation):
\begin{dashlist}
    \item \textbf{Selection} ($\sigma$): output a subset of \emph{rows} (horizontal filtering).
    \item \textbf{Projection} ($\pi$): output a subset of \emph{columns} (vertical filtering).
\end{dashlist}

\textbf{Binary operators} (two input relations):
\begin{dashlist}
    \item \textbf{Cross-product} ($\times$): concatenate every row of the first relation with every row of the second.
    \item \textbf{Set-difference} ($-$): output tuples in $R$ that are not in $S$.
    \item \textbf{Union} ($\cup$): output tuples in $R$ or in $S$ (or both).
\end{dashlist}

Since each operation returns a relation, \textbf{operations can be composed}.

\subsection{Selection ($\sigma$)}
Selection outputs only the rows that satisfy a given condition. The output schema is the same as the input.

\[ \sigma_{\text{condition}}(\textit{relation}) \]

\subsubsection*{Without condition}
The simplest expression $\sigma(\textit{instructor})$ returns all rows unchanged (equivalent to \lstinline[style=SQLStyle]{SELECT * FROM instructor}).

\subsubsection*{With a single condition}
\[ \sigma_{\textit{dept\_name} = \text{``Physics''}}(\textit{instructor}) \]

\begin{sql}
    SELECT * FROM instructor WHERE dept_name = 'Physics'
\end{sql}

\begin{dbtable}{}{clll}
\textbf{ID} & \textbf{name} & \textbf{dept\_name} & \textbf{salary} \\ \hline
22222 & Einstein & Physics & 95000 \\ \hline
33456 & Gold & Physics & 87000 \\
\end{dbtable}

\subsubsection*{With a compound condition}
Conditions can be combined with $\wedge$ (and), $\vee$ (or), and $\neg$ (not):

\[ \sigma_{\textit{dept\_name} = \text{``Physics''} \,\wedge\, \textit{salary} > 90000}(\textit{instructor}) \]

\begin{sql}
    SELECT * FROM instructor WHERE dept_name = 'Physics' AND salary > 90000
\end{sql}

\begin{dbtable}{}{clll}
\textbf{ID} & \textbf{name} & \textbf{dept\_name} & \textbf{salary} \\ \hline
22222 & Einstein & Physics & 95000 \\
\end{dbtable}

\subsection{Projection ($\pi$)}
Projection retains only the attributes in the projection list. The output schema contains exactly those attributes.

\[ \pi_{\text{attribute list}}(\textit{relation}) \]

\smallskip
\[ \pi_{\textit{ID, name, salary}}(\textit{instructor}) \]
corresponds to:
\begin{sql}
    SELECT ID, name, salary FROM instructor
\end{sql}
The result has three columns instead of four (\texttt{dept\_name} is dropped), but all 12 rows are kept because each row is still distinct.

\subsubsection*{Duplicate elimination}
In relational algebra, projection \textbf{eliminates duplicates} because the result must be a set of tuples.

\[ \pi_{\textit{dept\_name}}(\textit{instructor}) \]
reduces the 12 rows to 7 distinct department names. In SQL, duplicates are kept by default; to match the algebra, use \lstinline[style=SQLStyle]{SELECT DISTINCT}:
\begin{sql}
    SELECT DISTINCT dept_name FROM instructor
\end{sql}

\subsection{Composing operators}
Since every operator returns a relation, the output of one can feed into another. Selection followed by projection gives specific columns from filtered rows:

\[ \pi_{\textit{name}}\!\left(\sigma_{\textit{dept\_name} = \text{``Physics''}}(\textit{instructor})\right) \]

\begin{sql}
    SELECT name FROM instructor WHERE dept_name = 'Physics'
\end{sql}

\begin{dbtable}{}{l}
\textbf{name} \\ \hline
Einstein \\ \hline
Gold \\
\end{dbtable}

First $\sigma$ filters to Physics instructors, then $\pi$ keeps only the \texttt{name} column. The order matters: projecting first would discard the column needed for selection.

\newpage
\subsection{Rename operator ($\rho$)}
The rename operator ($\rho$) renames a relation or its attributes. Since algebraic expressions can yield unnamed relations natively, renaming is essential to avoid naming conflicts or to self-reference.

\begin{dashlist}
    \item Renames the list of attributes in the form $\text{oldname} \rightarrow \text{newname}$ or $\text{position} \rightarrow \text{newname}$.
    \item Can be used to rename the output relation itself entirely.
    \item \textbf{Output schema}: same as input except for the renamed attributes.
    \item \textbf{Output instances}: returns the exact same tuples as the input.
\end{dashlist}

Renaming the relation itself to $i$:
\[ \rho_{i}(\textit{instructor}) \]

Renaming the \textit{ID} attribute to \textit{instructor.ID}:
\[ \rho_{ID \rightarrow \textit{instructor.ID}}(\textit{instructor}) \]

In SQL, this is achieved using the \lstinline[style=SQLStyle]{AS} keyword in the \lstinline[style=SQLStyle]{SELECT} or \lstinline[style=SQLStyle]{FROM} clauses:
\begin{sql}
    SELECT i.name FROM instructor AS i WHERE i.ID = ...
\end{sql}

\subsection{Union operator ($\cup$)}
The union operator takes two input relations, which \textbf{must} be \emph{union-compatible}:
\begin{dashlist}
    \item They must have the same number of fields.
    \item The ``corresponding'' fields must have exactly the same type.
\end{dashlist}

Like projection, union in strict relational algebra implies \emph{duplicate elimination}.
\[ c1 \cup c2 \]

In SQL, \lstinline[style=SQLStyle]{UNION} implicitly eliminates duplicates, while \lstinline[style=SQLStyle]{UNION ALL} is faster but keeps them.
\begin{sql}
    (SELECT * FROM c1)
    UNION
    (SELECT * FROM c2)
\end{sql}

\subsection{Set-difference operator ($-$)}
The set-difference operator outputs the set of tuples in $c1$ which are not in $c2$.
\begin{dashlist}
    \item The two input relations must also be \emph{union-compatible}.
    \item Set difference is \textbf{not} commutative ($c1 - c2 \neq c2 - c1$).
\end{dashlist}

\[ c1 - c2 \]

In SQL, this corresponds to the \lstinline[style=SQLStyle]{EXCEPT} clause:
\begin{sql}
    (SELECT * FROM c1)
    EXCEPT
    (SELECT * FROM c2)
\end{sql}

\subsection{Compound operators}
Alongside the five basic operators ($\sigma, \pi, \times, -, \cup$), there are several additional \textbf{compound operators}. These add \textbf{no computational power} to the language (they can all be expressed using the basic operators), but they serve as useful, compact shorthands.

\subsubsection*{Intersection ($\cap$)}
Intersection outputs the overlapping tuples from two \emph{union-compatible} relations.
\[ R \cap S = R - (R - S) \]

\begin{sql}
    (SELECT * FROM c1)
    INTERSECT
    (SELECT * FROM c2)
\end{sql}

\subsection{Cross-product operator ($\times$)}
The Cartesian product $S \times R$ pairs every single row of $S$ with every single row of $R$.
\begin{dashlist}
    \item \textbf{Result schema}: has one field per field of $S$ and $R$, inheriting the field names linearly.
    \item \textbf{Naming conflicts}: Both $S$ and $R$ may have fields sharing the exact same name. In this case, we use the renaming operator ($\rho$) to disambiguate.
\end{dashlist}

\[ \textit{instructor} \times \textit{teaches} \]

\subsection{Join operator ($\bowtie$)}
A simple $\textit{instructor} \times \textit{teaches}$ associates \emph{every} instructor tuple with \emph{every} teaches tuple. Most of these rows hold information about an instructor paired with a course they never taught!

To get only tuples of instructors and the specific courses they actually taught, we apply a selection:
\[ \sigma_{\textit{instructor.id} = \textit{teaches.id}} (\textit{instructor} \times \textit{teaches}) \]

Because this pattern ($\sigma$ after $\times$) is overwhelmingly common in databases, it gets its own dedicated operator called a \textbf{Join}.
\[ \textit{instructor} \bowtie_{\textit{instructor.id} = \textit{teaches.id}} \textit{teaches} \]

If the matched columns have the identical name in both tables, the condition can be completely omitted. This is called a \textbf{Natural Join}.

In SQL, a join is explicitly specified using \lstinline[style=SQLStyle]{JOIN}:
\begin{sql}
    SELECT * FROM instructor
    JOIN teaches ON instructor.ID = teaches.ID
\end{sql}

\subsection{Complex queries}
Operators can be chained to answer sophisticated questions. Consider: ``Find the names of all instructors in the Music department together with the course title of all the courses that they teach.''

This query involves three relations (\texttt{instructor}, \texttt{teaches}, \texttt{course}) and requires joining them before projecting the desired columns.

Two equivalent formulations:

\begin{enumerate}[label=(\alph*)]
    \item Apply the selection \emph{after} the full join:
    \[ \pi_{\textit{name, title}}\!\left(\sigma_{\textit{dept\_name} = \text{``Music''}}\!\left(\textit{instructor} \bowtie (\textit{teaches} \bowtie \pi_{\textit{course\_id, title}}(\textit{course}))\right)\right) \]
    \item Push the selection \emph{before} the join:
    \[ \pi_{\textit{name, title}}\!\left(\left(\sigma_{\textit{dept\_name} = \text{``Music''}}(\textit{instructor})\right) \bowtie \left(\textit{teaches} \bowtie \pi_{\textit{course\_id, title}}(\textit{course})\right)\right) \]
\end{enumerate}

Both return the same result, but formulation (b) is significantly more efficient: the selection filters the instructor relation \emph{before} the expensive join, drastically reducing the number of intermediate tuples.

\subsubsection*{Query optimization}
This illustrates a key principle: \textbf{query optimization determines performance}. The DBMS query optimizer automatically rewrites expressions (e.g.\ pushing selections down) to find an efficient execution plan. The two expression trees below show the difference visually:

\begin{minipage}[t]{0.48\textwidth}
\centering
\input{\subfix{schemas/tree_initial.tex}}

\smallskip
(a) Initial: $\sigma$ applied \emph{after} joining all three relations — processes every tuple.
\end{minipage}
\hfill
\begin{minipage}[t]{0.48\textwidth}
\centering
\begin{tikzpicture}[
    node distance=1.1cm and 0.8cm,
    every node/.style={font=\small},
    leaf/.style={font=\small\itshape},
    >=Latex
]
    % Nodes top-down
    \node (pi) {$\pi_{\textit{name, title}}$};
    \node[below=of pi] (join1) {$\bowtie$};
    \node[below left=of join1, draw, ellipse, red, thick, inner sep=2pt] (sigma) {$\sigma_{\textit{dept\_name}=\text{``Music''}}$};
    \node[below right=of join1] (join2) {$\bowtie$};
    \node[leaf, below=of sigma] (inst) {instructor};
    \node[leaf, below left=of join2] (teach) {teaches};
    \node[below right=of join2] (picourse) {$\pi_{\textit{course\_id, title}}$};
    \node[leaf, below=of picourse] (course) {course};

    % Edges
    \draw (pi) -- (join1);
    \draw (join1) -- (sigma);
    \draw (join1) -- (join2);
    \draw (sigma) -- (inst);
    \draw (join2) -- (teach);
    \draw (join2) -- (picourse);
    \draw (picourse) -- (course);
\end{tikzpicture}
}

\smallskip
(b) Optimized: $\sigma$ pushed \emph{down} to filter \textit{instructor} before the join — far fewer tuples flow upward.
\end{minipage}

\medskip
Both trees compute the same result, but (b) is far more efficient because the expensive join operates on a tiny filtered relation instead of the full table.

\newpage
\subsection{Division operator ($/$)}
Division is a compound operator useful for expressing ``for all'' queries (e.g., ``find all suppliers who supply \emph{every} part'').

\begin{concept}{Division ($A / B$)}
Given relations $A(x, y)$ and $B(y)$, the division $A / B$ returns all values of $x$ such that for \textbf{every} tuple $b$ in $B$, the tuple $(x, b)$ exists in $A$.
\end{concept}

\subsubsection*{Worked example}
Consider relation $A(\texttt{sno}, \texttt{pno})$ and three divisors $B_1 \subseteq B_2 \subseteq B_3$:

{\small
\begin{minipage}[t]{0.22\textwidth}
\begin{dbtable}{$A$}{ll}
\textbf{sno} & \textbf{pno} \\ \hline
s1 & p1 \\ \hline
s1 & p2 \\ \hline
s1 & p3 \\ \hline
s1 & p4 \\ \hline
s2 & p1 \\ \hline
s2 & p2 \\ \hline
s3 & p2 \\ \hline
s4 & p2 \\ \hline
s4 & p4 \\
\end{dbtable}
\end{minipage}
\hfill
\begin{minipage}[t]{0.74\textwidth}
\renewcommand{\arraystretch}{1.2}
\begin{center}
\begin{tabular}{@{}lll@{}}
$B_1 = \{\text{p2}\}$ & $B_2 = \{\text{p2, p4}\}$ & $B_3 = \{\text{p1, p2, p4}\}$ \\[4pt]
$A / B_1 = \{\text{s1, s2, s3, s4}\}$ & $A / B_2 = \{\text{s1, s4}\}$ & $A / B_3 = \{\text{s1}\}$
\end{tabular}
\end{center}

\smallskip
As $B$ grows (more conditions to satisfy), the result of $A/B$ shrinks: only suppliers matching \emph{every} part in $B$ survive.
\end{minipage}
}

\subsubsection*{Expressing division with basic operators}
Division can be expressed using only $\pi$, $\times$, and $-$.

\textbf{Idea}: find all $x$ values that are \emph{not} ``disqualified''. An $x$ is disqualified if pairing it with some $y \in B$ yields a tuple \emph{not} in $A$.

Disqualified $x$ values:
\[ \pi_x\!\left((\pi_x(A) \times B) - A\right) \]

Final result:
\[ A / B \;=\; \pi_x(A) \;-\; \pi_x\!\left((\pi_x(A) \times B) - A\right) \]

\subsubsection*{Step by step with $B_3 = \{\text{p1, p2, p4}\}$}
\begin{dashlist}
    \item $\pi_{\textit{sno}}(A) = \{\text{s1, s2, s3, s4}\}$
    \item $\pi_{\textit{sno}}(A) \times B_3$ produces all $4 \times 3 = 12$ possible (sno, pno) pairs.
    \item $(\pi_{\textit{sno}}(A) \times B_3) - A$ keeps only the \emph{missing} pairs: $(\text{s2, p4}),\, (\text{s3, p1}),\, (\text{s3, p4}),\, (\text{s4, p1})$.
    \item Projecting on \texttt{sno} gives the disqualified suppliers: $\{\text{s2, s3, s4}\}$.
    \item $\pi_{\textit{sno}}(A) - \{\text{s2, s3, s4}\} = \{\text{s1}\}$.
\end{dashlist}

\end{document}
